\problem{Émission de rayonnement par un corps macroscopique}

\question{Calculez le nombre de photons $dN(\omega)$, dont la
fréquence est $\omega$ (à $d\omega$ près), qui sortent à travers du trou dans une intervalle de temps $dt$.}

Nous savons que la densité en fonction de $\mathbf{k}$ est :
\[
    \frac{dN(\mathbf{k})}{V} &= \sum_{s} \rho(\mathbf{k})\bar{n}_{\mathbf{k}}^B d^3k. \notag
\]

Or, nous savons que $\rho = \frac{V}{(2\pi )^3}$ et que la somme des polarisations est 2. 
Nous avons donc :
\[
    dN(\mathbf{k}) = \bar{n}_{\mathbf{k}}^B \frac{2}{(2\pi )^3} d^3k. \notag
\]

Maintenant, déterminons le nombre de photons qui sort par le trou :
\[
    dN(\mathbf{k}) &= \bar{n}_{\mathbf{k}}^B \frac{2}{(2\pi )^3} d^3k c \, dt \, dS \, \cos(\theta) \notag \\
    &= \int_{0}^{\frac{\pi}{2}} \sin(\theta) \cos(\theta) \, d\theta \int_{0}^{2\pi} d\phi \bar{n}_k^B \frac{2}{(2\pi )^3} c k^2 \, dk \, dt \, dS. \notag
\]

Avec Wolfram Alpha, nous trouvons que le résultat des intégrales est de $\pi$. Avec cela, nous avons :
\[
    dN(\mathbf{k}) &= \pi \bar{n}_k^B \frac{2}{(2\pi )^3}c k^2 \, dk \, dt \, dS \notag \\
    &= \frac{c \, dS \, dt}{4\pi^2} \bar{n}_k^B k^2 \, dk. \notag
\]

Avec la relation $\omega = ck$ et $d\omega = c \, dk$, nous obtenons :
\[
    dN(\omega) &= \frac{c \, dS \, dt}{4\pi^2} \bar{n}_{\omega}^B \frac{\omega^2}{c^3} \, d\omega \notag \\
    &= \frac{\omega^2}{4\pi^2 c^2}\frac{1}{e^{\beta \hbar \omega} - 1} \, d\omega \, dt \, dS.
\]

\question{Quelle est l'énergie portée par les photons de (1) ?}

Puisque l'énergie d'un photon est de $\hbar \omega$, l'énergie totale est :
\[ 
    dE(\omega) &= \hbar \omega \, dN(\omega) \notag \\
    &= \frac{\hbar \omega^3}{4\pi^2 c^2}\frac{1}{e^{\beta \hbar \omega} - 1} \, d\omega \, dt \, dS.
\]

\question{À partir de (2), calculez la puissance totale (incluant toutes les fréquences) par unité
d'aire émise à travers du trou}

La puissance par unité d'aire est l'énergie total divisé par le temps ainsi que par l'aire. 
\[
    P(\omega) &= \frac{dE}{dt \, dS} \notag \\
    &= \frac{\hbar \omega^3}{4\pi^2 c^2}\frac{1}{e^{\beta \hbar \omega} - 1} \, d\omega. \notag
\]

En intégrant sur les fréquences, nous obtenons la puissance totale :
\[
    P_{tot} &= \frac{\hbar}{4\pi^2 c^2} \int_{0}^{\infty} \frac{\omega^3}{e^{\beta \hbar \omega} - 1} \, d\omega \notag \\
    &= \frac{\hbar}{4\pi^2 c^2} \qty(\frac{k_B T}{\hbar})^4 \int_{0}^{\infty} \frac{x^3}{e^x - 1} \, dx \notag \\
    &= \frac{k_B^4 T^4}{4\pi^2 c^2 \hbar^3}\frac{\pi^4}{15} \notag \\
    &= \frac{\pi^2 k_B^4T^4}{60 c^2 \hbar^3}.
\]

\question{ À partir de (3), montrez que la puissance rayonnée par un corps sphérique de rayon R.}

La puissance totale émise par unité d'aire est :
\[
    P = \sigma T^4, \notag
\]

avec $\sigma$ étant la constante de Stefan-Boltzmann. Pour calculer la puissance totale d'une sphère, 
Il faut multiplier cette puissance par l'aire d'une sphère, soit $4\pi R^2$. Nous avons donc :
\[
    P = 4\pi R^2 \sigma T^4.
\]
\clearpage

